\pagestyle{plain}
\usepackage[a4paper,hscale=0.77,vscale=0.84,tmargin=2cm,bmargin=2cm]{geometry}
\usepackage{natbib}
\usepackage{booktabs}
\usepackage{paralist}
\usepackage{multicol}

\newcommand{\ConRef}{
  \nocite{Strang:18.06}
  \bibliography{referencias/referencias}
  \bibliographystyle{plainnat}
}

\newcommand{\eng}[2]{{#1}}

\newcommand{\corte}{ 
  \begin{flushright}
    \emph{\eng
      {El examen continúa en la siguiente página}
      {The exam continues in the next page}}
      $\rightarrow$
  \end{flushright}
  \pagebreak}

\usepackage[spanish]{babel}
\deactivatetilden
\decimalpoint 
\usepackage[spanish]{varioref}
\renewcommand\shorthandsspanish{}

\usepackage[pdftex,nodirectory,latextoc,spanish]{acrotex/web}

%%%%%%%%%%%%%%%%%%%%%%%%%%%%%%%%%%%%%%%%%%%%%%%%%%%%%%%%%%%%%% 
%% UNA Y SOLO UNA DE ESTAS TRES LÍNEAS DEBE ESTAR DESCOMENTADA
%
% \usepackage[pdftex,forpaper,nohiddensolutions,solutionsafter,spanish]{exerquiz}   % soluciones a continuación de los enunciados
% \usepackage[pdftex,forpaper,nosolutions,spanish]{exerquiz}                        % sin soluciones (versión de examen para reprografía)
\usepackage[pdftex,forpaper,spanish]{acrotex/exerquiz}                                    % soluciones al final del documento
%%%%%%%%%%%%%%%%%%%%%%%%%%%%%%%%%%%%%%%%%%%%%%%%%%%%%%%%%%%%%%

%%%%%%%%%%%%%%%%%%%%%%%%%%%%%%%%%%%%%%%%%%%%%%%%%%%%%%%%%%%%%
%% UNA Y SOLO UNA DE ESTAS DOS LÍNEAS DEBE ESTAR DESCOMENTADA
\newenvironment{procedencia}[1]{\noindent }                                         % se oculta la procedencia del ejercicio
% \newenvironment{procedencia}[1]{\noindent \emph{#1}}                              % se indica la procedencia si está puesta
%%%%%%%%%%%%%%%%%%%%%%%%%%%%%%%%%%%%%%%%%%%%%%%%%%%%%%%%%%%%%

\newenvironment{pista}{\textbf{Pista:}}{}

\newcommand{\verdaderofalso}{Verdadero o falso
  (\emph{solo puntuan respuestas correctamente justificadas; una respuesta sin explicación tendrá una calificación de cero puntos}).
  \smallskip}


\newcommand\myquote[2]{%
  \smallskip
  
  \par
  \hspace{-0.9cm}
  \begin{minipage}{\textwidth-1.5\parindent\relax}
    \rule{\textwidth}{.2pt}
    
    #1\vspace{-1.2ex} 
    \rule{\textwidth}{.2pt}
    \par\vspace{0.5ex}
  \end{minipage}
  \par}

%%%%%%%%%%%%%%%%%%%%%%%%%%%%%%%%%%%%%%%%%%%%%%%%%%%%%%%%%%%%%%%%%%
%%%% Indicador de puntos por pregunta y Contador de puntos totales
\usepackage{xfp} % Para cálculos con decimales

% Configuración del acumulador global para puntos
\ExplSyntaxOn
\fp_new:N \g_puntos_total_fp % Crea un acumulador global para puntos decimales

% Redefinimos el comando \puntos para que haga la cuenta de puntos acumulados
\NewDocumentCommand{\puntos}{m}{
  \PTsHook{($\eqPTs^{\text{pts}}$)}\PTs{#1} % Tu comando original
  \fp_gadd:Nn \g_puntos_total_fp {#1} % Suma los puntos al acumulador global
}

% Comando para mostrar la suma total de puntos
\NewDocumentCommand{\totalpuntos}{}{Puntos totales en el examen:\; \fp_use:N \g_puntos_total_fp}
\ExplSyntaxOff
%%%%%%%%%%%%%%%%%%%%%%%%%%%%%%%%%%%%%%%%%%%%%%%%%%%%%%%%%%%%%%%%%%

%%% PARA LA SECCIÓN DE CALIFICACIÓN DE LA CABECERA
\usepackage{totcount}
\def\computemycount{}
\expandafter\ifx\csname computemycount\endcsname\relax
\usetotcountfile{mycount.aux}
\newcounter{mycount} % so that increments of this counter don't break
\else
\newtotcounter[auxfile=mycount.aux]{mycount}
\fi
\regtotcounter{eqexno}

\newcounter{ultimo}
\newcounter{EjNum}
\setcounter{EjNum}{1}
\usepackage{pgffor}
\newcommand{\contador}[1]{%
  \def\temp{}%
  \foreach \i in {1,...,#1}
    {%
      \expandafter\gdef\expandafter\temp\expandafter{\temp\noindent\\[5pt]\textbf{\arabic{EjNum}.-} \stepcounter{EjNum}}
    }%
  \temp}

\usepackage{mathtools}
\usepackage{systeme}

\usepackage{nacal/nacal}

\usepackage{tikz}
\usetikzlibrary{arrows,positioning}
\usepackage{tikz-3dplot}
\usepackage{gnuplot-lua-tikz}
%\pgfplotsset{compat=1.10}
\usetikzlibrary{intersections}
\usetikzlibrary{calc}
\usepackage{tkz-euclide}

\usepackage[autoprint=false]{pythontex}
\pythontexcustomc{pyconsole}{from nacal import *}
\pythontexcustomc{py}{from nacal import *}
\pythontexcustomc{sympy}{from nacal import *}

\newcommand{\PreguntasLargas}{
  \setcounter{eqexno}{0}
  \renewcommand\exlabel{Ejercicio}
}

\newcommand{\PreguntasCortas}{
  \vspace{.5cm}
  
  Conteste a TODOS los apartados de TODOS los \framebox{\textsl{CONJUNTOS DE PREGUNTAS CORTAS}}

  \setcounter{eqexno}{0}
  \renewcommand\exlabel{Conjunto de preguntas}
  \vspace{.5cm}
}

\endinput
